% Technical note: how much of the headline carbon saving is real?
% Compile: pdflatex (MiKTeX at AppData\Local\Programs\MiKTeX\miktex\bin\x64\pdflatex.exe)
\documentclass[11pt,a4paper]{article}
\usepackage[margin=2.4cm]{geometry}
\usepackage[T1]{fontenc}
\usepackage{lmodern}
\usepackage{titlesec}
\usepackage{enumitem}
\usepackage{xcolor}
\usepackage{booktabs}
\usepackage[hidelinks]{hyperref}
\usepackage{parskip}

\definecolor{rulecol}{RGB}{49,86,127}
\titleformat{\section}{\large\bfseries\color{rulecol}}{\thesection.}{0.5em}{}
\titlespacing{\section}{0pt}{1.4em}{0.4em}
\titleformat{\subsection}{\normalsize\bfseries}{}{0pt}{}
\titlespacing{\subsection}{0pt}{1em}{0.2em}
\setlist[itemize]{leftmargin=1.3em, itemsep=0.15em, topsep=0.3em}
\pagestyle{plain}

\begin{document}

\begin{center}
  {\LARGE\bfseries How much of a 44\% carbon saving is real?}\\[6pt]
  {\large Uncertainty in a low-carbon concrete mix screening tool}\\[10pt]
  {\normalsize Harvey Sohal \quad\textbullet\quad MEng Civil Engineering, Imperial College London}\\[2pt]
  {\small \href{mailto:harvey@sohal.info}{harvey@sohal.info} \quad\textbullet\quad
   \href{https://github.com/harvey-a11y/carbon-mix-optimiser}{github.com/harvey-a11y/carbon-mix-optimiser}}\\[2pt]
  {\small August 2026}
\end{center}

\vspace{0.6em}
\hrule
\vspace{0.8em}

\textbf{Summary.} I built a tool that searches concrete mixes for the lowest embodied carbon
that still meets a strength class and a durability requirement. For a C32/40 XC3/4 specification
it found a mix saving 43.9\% against a CEM~I baseline. This note is about whether that number
deserves a decimal place. It does not. Running the same search a few hundred times over
plausible ranges for the uncertain inputs gives a median saving of about 41\% and a 90\%
interval of roughly 8\% to 62\%. More usefully, almost all of that spread comes from one place:
the strength model, not the carbon data. Refining the carbon factors would move the answer by
about five percentage points. Calibrating the strength model would move it by about fifty-five.

\section{What the tool does}

Given a target strength class and an exposure class, it enumerates candidate mixes over total
binder content, free water/cement ratio, and GGBS and fly ash replacement fractions. For each
one it estimates 28-day strength using an Abrams-type relationship with the EN~206 $k$-value
weighting for supplementary cementitious materials, screens it against BS~8500-style limits on
effective w/c and minimum binder content, and computes embodied carbon from a mass balance over
one cubic metre.

Feasible mixes are ranked by carbon. The tool reports the best one and its saving against the
lowest-carbon feasible mix that uses CEM~I only.

For C32/40 in exposure class XC3/XC4 it evaluates 57{,}040 combinations and returns:

\begin{center}
\begin{tabular}{lrr}
\toprule
 & \textbf{CEM I baseline} & \textbf{Best mix} \\
\midrule
Total binder (kg/m\textsuperscript{3}) & 300 & 300 \\
GGBS (\% of binder)                    & 0   & 50 \\
Free w/c                               & 0.43 & 0.35 \\
Embodied carbon (kgCO\textsubscript{2}e/m\textsuperscript{3}) & 283.1 & 158.8 \\
\midrule
\multicolumn{2}{l}{\textbf{Saving}} & \textbf{43.9\%} \\
\bottomrule
\end{tabular}
\end{center}

\section{Why I stopped trusting that number}

Two of the inputs are, by my own documentation, not calibrated to anything.

The strength model uses $f_{cm} = K_1 / K_2^{\,w/c}$ with $K_1 = 110$ and $K_2 = 10$. Those
constants give sensible-looking answers for a CEM~I concrete and that is the only justification
for them. Real strength depends on cement source and strength class, aggregate, admixtures,
curing and age, none of which the model sees.

The carbon factors are order-of-magnitude figures in the style of the ICE database. Cement is
around 0.91 kgCO\textsubscript{2}e/kg and GGBS around 0.08, but published values vary with plant,
fuel mix, transport and, for GGBS especially, with how you allocate the emissions of steelmaking
between the steel and the slag.

Quoting 43.9\% on top of that is claiming three significant figures from inputs that do not
support two. So I made the tool report a range instead.

\section{Method}

I sample the six carbon factors and the two Abrams constants from independent uniform ranges,
re-run the entire grid search for each draw, and collect the resulting saving.

Uniform rather than normal is deliberate. These are ranges of plausible values, not measurement
errors around a known mean, and assuming a shape would add a second unjustified assumption on
top of the first. Independence is also an assumption and a questionable one, since cement and
GGBS share a supply chain and a fuel mix. Both choices are recorded in the code.

\begin{center}
\begin{tabular}{lrr}
\toprule
\textbf{Input} & \textbf{Low} & \textbf{High} \\
\midrule
CEM I (kgCO\textsubscript{2}e/kg)   & 0.75  & 0.95 \\
GGBS                                & 0.05  & 0.13 \\
Fly ash                             & 0.004 & 0.030 \\
Aggregate, sand                     & 0.003 & 0.008 \\
$K_1$ (MPa)                         & 95    & 130 \\
$K_2$                               & 8     & 13 \\
\bottomrule
\end{tabular}
\end{center}

\section{Results}

Four hundred draws for C32/40, XC3/XC4. Three of them found no feasible blended mix at all and
are excluded.

\begin{center}
\begin{tabular}{lr}
\toprule
Median saving        & 41.0\% \\
Mean saving          & 39.8\% \\
90\% interval        & 7.6\% to 62.3\% \\
Full range observed  & 0.0\% to 64.8\% \\
\bottomrule
\end{tabular}
\end{center}

The interval is very wide. A tool that can return anything between no saving at all and roughly
two thirds, depending on constants nobody has measured, is not yet an instrument for making
decisions. That is worth saying plainly, because the deterministic answer looks a great deal
more confident than that.

\subsection{Which input is responsible}

The useful question is not how wide the interval is but which assumption makes it wide. Varying
one group at a time, holding the other at its default, over 400 draws each:

\begin{center}
\begin{tabular}{lcr}
\toprule
\textbf{Varying} & \textbf{90\% interval} & \textbf{Width} \\
\midrule
Carbon factors only       & 40.2\% to 44.8\% & 4.6 points \\
Strength calibration only & \phantom{0}7.4\% to 63.0\% & 55.6 points \\
Both                      & \phantom{0}8.6\% to 61.9\% & 53.4 points \\
\bottomrule
\end{tabular}
\end{center}

This is the result that changed how I think about the tool. The carbon factors, which are the
part everyone argues about and the part I expected to dominate, contribute about five percentage
points of spread. The uncalibrated strength model contributes about fifty-five, and adding the
carbon uncertainty on top of it barely widens the interval further.

The mechanism is straightforward once you see it. The strength model does not just shift the
answer, it changes which mixes are feasible at all. A slightly optimistic calibration lets high
replacement levels reach the target strength and the saving is large; a slightly pessimistic one
rules them out and the best feasible mix is barely better than CEM~I. The carbon factors only
rescale a set of options that has already been chosen.

\subsection{What that implies for the work}

Improving the carbon data would be close to wasted effort while $K_1$ and $K_2$ remain
uncalibrated. The next useful thing to do with this tool is to fit the strength model to real
28-day cube data for a small number of mixes, and only then worry about which ICE revision the
factors came from. That is not what I would have guessed before running it.

\section{Two other things the tool was ignoring}

\subsection{Transport}

The headline is cradle-to-gate, modules A1--A3, which is the standard basis for comparing mixes.
It has a specific blind spot here: GGBS usually travels further than cement in the UK, because
supply is concentrated around a small number of works and import terminals while cement plants
are more widely distributed. A tool that recommends 50\% GGBS and never mentions haulage is
answering an easier question than the one that matters on site.

Module A4 is now available as an option, defaulting to zero so the A1--A3 figure is unchanged
unless distances are supplied. At an indicative 0.11 kgCO\textsubscript{2}e per tonne-kilometre:

\begin{center}
\begin{tabular}{lcc}
\toprule
\textbf{Scenario} & \textbf{A1--A3} & \textbf{A1--A4} \\
\midrule
Both materials local (80 km)            & 43.9\% & 42.1\% \\
Cement 80 km, GGBS 300 km               & 43.9\% & 40.9\% \\
Cement 80 km, GGBS 500 km               & 43.9\% & 39.8\% \\
\bottomrule
\end{tabular}
\end{center}

Transport erodes the saving but does not overturn it, which is a smaller effect than I expected
and worth knowing either way. It is also much smaller than the strength-calibration uncertainty
above, which puts it in perspective.

\subsection{Early-age strength}

Every calculation in the tool is a 28-day calculation. High-GGBS mixes gain strength slowly, and
that governs formwork striking, early loading and cold-weather working. The recommended mix here
is 50\% GGBS, for which an indicative 7-day strength is around 50\% of the 28-day value against
roughly 65\% for CEM~I.

The tool now says so whenever it recommends a mix above 50\% replacement. It deliberately does
not offer a striking time, because nothing in it is calibrated well enough to justify one. That
needs cube or maturity data from the actual mix at the actual temperature.

\section{Verification}

The Python package could not be executed on the machine this was written on, because Python is
blocked there by an application control policy. Rather than quote figures nobody had computed, I
re-implemented the model independently in JavaScript, working from the module docstrings rather
than translating the code line by line, and used that to produce every number in this note.

The second implementation reproduces the deterministic 43.9\% exactly. It is committed as
\texttt{tools/verify\_model.js} so the two can be compared. Continuous integration runs the
Python test suite, 83 tests, on Python 3.11, 3.12 and 3.13.

\section{What this is not}

This is a screening tool and a piece of coursework, not a design tool. It does not model
cover, cement designations, intended working life, aggregate size, admixtures, workability,
pumpability, shrinkage, creep or carbonation. Its durability limits are BS~8500-\emph{style}
rather than a transcription of the standard. No number in it should be used for a real mix
without trial mixes and a chartered engineer.

The honest position after doing this work is that the tool is good at asking which direction to
go and bad at saying how far. That is still useful, as long as it is the claim being made.

\vspace{1em}
\hrule
\vspace{0.6em}
{\small All figures reproducible: \texttt{node tools/verify\_model.js}, or
\texttt{python -m carbonmix -{}-class C32/40 -{}-exposure XC3\_XC4 -{}-sensitivity 400}.}

\end{document}
